% AST Compiler + Duck Types + Čech Algo + E2E Trace (Ch 37-40)
% Lines 7829-8659 of main.tex
% This file is for agent editing — will be merged back

\chapter{The AST Compiler in Detail}
\label{ch:compiler-detail}

This chapter specifies the compilation from Python AST to Z3 $\PyObj$
terms with enough detail that a reader could implement it from scratch.

\section{Compilation State}

The compiler maintains an \emph{environment} --- a mapping from
variable names to Z3 terms:

\begin{definition}[Compilation Environment]
\[
  \mathsf{Env} = \{T : \mathsf{Theory},\;
    \mathsf{bindings} : \mathsf{str} \to \PyObj,\;
    \mathsf{func\_name} : \mathsf{str},\;
    \mathsf{is\_rec} : \Bool\}
\]
\end{definition}

The environment is \emph{immutable} (in the cubical sense): branching
creates copies.  When an \texttt{if} statement has both a \texttt{then}
and \texttt{else} branch, the compiler copies the environment for each
branch and merges the results with $\mathsf{If}$ expressions.

\section{Expression Compilation Rules}

Each Python expression AST node maps to a Z3 $\PyObj$ term.  We give
the complete rule set:

\subsection{Rule E1: Literals}
\[
  \den{n}_\mathsf{env} = \mathsf{IntObj}(n) \quad
  \den{s}_\mathsf{env} = \mathsf{StrObj}(s) \quad
  \den{\mathsf{True}}_\mathsf{env} = \mathsf{BoolObj}(\top) \quad
  \den{\mathsf{None}}_\mathsf{env} = \mathsf{NoneObj}
\]

\subsection{Rule E2: Variables}
\[
  \den{x}_\mathsf{env} = \begin{cases}
    \mathsf{env.bindings}[x] & \text{if $x \in \mathsf{env.bindings}$} \\
    \mathsf{FreshConst}(\PyObj) & \text{otherwise (undeclared variable)}
  \end{cases}
\]

\subsection{Rule E3: Binary Operations}
\[
  \den{a \oplus b}_\mathsf{env} = \mathsf{binop}(\llbracket \oplus \rrbracket,
    \den{a}_\mathsf{env}, \den{b}_\mathsf{env})
\]
where $\llbracket \oplus \rrbracket$ maps Python operators to integer
operation tags:
\begin{center}
\small
\begin{tabular}{l r l r l r l r}
  \texttt{+} & 1 & \texttt{-} & 2 & \texttt{*} & 3 &
  \texttt{//} & 5 \\
  \texttt{\%} & 6 & \texttt{<} & 20 & \texttt{<=} & 21 &
  \texttt{>} & 22 \\
  \texttt{>=} & 23 & \texttt{==} & 24 & \texttt{!=} & 25 &
  \texttt{|} & 12 \\
\end{tabular}
\end{center}

\subsection{Rule E4: Unary Operations}
\[
  \den{-a}_\mathsf{env} = \mathsf{unop}(\mathsf{NEG}, \den{a}_\mathsf{env})
  \qquad
  \den{\mathsf{not}\; a}_\mathsf{env} = \mathsf{unop}(\mathsf{NOT},
    \den{a}_\mathsf{env})
\]

\subsection{Rule E5: Comparisons (Chained)}
For $a < b \leq c$:
\[
  \den{a < b \leq c}_\mathsf{env} = \mathsf{If}(
    \mathsf{truthy}(\mathsf{binop}(\mathsf{LT}, \den{a}, \den{b})),\;
    \mathsf{binop}(\mathsf{LE}, \den{b}, \den{c}),\;
    \mathsf{BoolObj}(\bot))
\]

\subsection{Rule E6: Boolean Operations}
\begin{align}
  \den{a \;\mathsf{and}\; b}_\mathsf{env} &= \mathsf{If}(
    \mathsf{truthy}(\den{a}), \den{b}, \den{a}) \\
  \den{a \;\mathsf{or}\; b}_\mathsf{env} &= \mathsf{If}(
    \mathsf{truthy}(\den{a}), \den{a}, \den{b})
\end{align}
Note: Python's \texttt{and}/\texttt{or} return the \emph{value}, not
a boolean.  The Z3 encoding faithfully preserves this.

\subsection{Rule E7: Conditional Expression}
\[
  \den{t \;\mathsf{if}\; c \;\mathsf{else}\; f}_\mathsf{env} =
    \mathsf{If}(\mathsf{truthy}(\den{c}), \den{t}, \den{f})
\]

\subsection{Rule E8: Function Calls}

Function calls have the most complex compilation:

\begin{enumerate}
  \item \textbf{Self-recursive call}: If $\mathsf{env.is\_rec}$ and the
    callee is $\mathsf{env.func\_name}$, create a RecFunction
    application: $\mathsf{rec\_f}(\den{a_1}, \ldots, \den{a_n})$.

  \item \textbf{Nested function call}: If the callee is in the
    environment as $(\mathsf{\_\_func\_\_}, \mathsf{AST})$, invoke
    $\mathsf{inline\_call}$: substitute arguments into the body and
    compile the body.

  \item \textbf{isinstance}: If the callee is \texttt{isinstance} with
    two arguments, compile to:
    $\mathsf{BoolObj}(\mathsf{tag}(\den{x}) = \tau)$ where $\tau$
    is the type tag corresponding to the second argument.

  \item \textbf{Shared builtins}: If the callee name is in
    $\mathsf{SHARED\_BUILTINS}$ (\texttt{sorted}, \texttt{len},
    \texttt{sum}, etc.), use the shared function symbol:
    $T.\mathsf{shared\_fn}(\mathsf{name}, n)(\den{a_1}, \ldots,
    \den{a_n})$.

  \item \textbf{Known builtins with Z3 semantics}: \texttt{abs},
    \texttt{int}, \texttt{bool}, \texttt{len}, \texttt{max},
    \texttt{min} compile to $\mathsf{unop}$ or $\mathsf{binop}$
    with the appropriate operation tag.

  \item \textbf{Other named calls}: Use a shared function symbol
    keyed by the callee name: $T.\mathsf{shared\_fn}(\mathsf{call\_name},
    n)(\den{a_1}, \ldots, \den{a_n})$.

  \item \textbf{Method calls}: Compile receiver and arguments, use
    shared method symbol:
    $T.\mathsf{shared\_fn}(\mathsf{meth\_name}, n+1)(\den{\mathsf{recv}},
    \den{a_1}, \ldots, \den{a_n})$.
\end{enumerate}

\subsection{Rule E9: Subscript}
\[
  \den{a[b]}_\mathsf{env} = \mathsf{binop}(\mathsf{GETITEM},
    \den{a}, \den{b})
\]

\subsection{Rule E10: Tuple/List Display}
\[
  \den{[a, b, c]}_\mathsf{env} = \mathsf{Pair}(\den{a},
    \mathsf{Pair}(\den{b}, \mathsf{Pair}(\den{c}, \mathsf{NoneObj})))
\]

\subsection{Rule E11: Comprehension}
\[
  \den{[e \;\mathsf{for}\; x \;\mathsf{in}\; \mathsf{xs}]}_\mathsf{env}
    = T.\mathsf{shared\_fn}(\mathsf{comp}_{h(e)}, 1)(\den{\mathsf{xs}})
\]
where $h(e) = \mathsf{hash}(\mathsf{ast.dump}(e)) \mod 2^{31}$ is
the canonical hash of the element expression.

\subsection{Rule E12: Attribute Access}
\[
  \den{a.\mathsf{attr}}_\mathsf{env} = T.\mathsf{shared\_fn}(
    \mathsf{attr\_attr}, 1)(\den{a})
\]


\section{Statement Compilation Rules}

Statements modify the environment and produce \emph{sections} (for
\texttt{return} statements).

\subsection{Rule S1: Assignment}
\[
  \mathsf{exec}(x = e, \mathsf{env}) = \mathsf{env}[x \mapsto \den{e}]
\]

For tuple unpacking $x, y = e$:
\[
  \mathsf{exec}((x, y) = e, \mathsf{env}) = \mathsf{env}[
    x \mapsto \mathsf{binop}(\mathsf{GETITEM}, \den{e},
      \mathsf{IntObj}(0)),\;
    y \mapsto \mathsf{binop}(\mathsf{GETITEM}, \den{e},
      \mathsf{IntObj}(1))]
\]

\subsection{Rule S2: Augmented Assignment}
\[
  \mathsf{exec}(x \mathrel{\oplus}= e, \mathsf{env}) = \mathsf{env}[
    x \mapsto \mathsf{binop}(\llbracket \oplus \rrbracket,
    \mathsf{env}[x], \den{e})]
\]

\subsection{Rule S3: Return}
\[
  \mathsf{exec}(\mathsf{return}\; e, \mathsf{env}, \mathsf{guard}) =
    \mathsf{Section}(\mathsf{guard}, \den{e})
\]
This terminates the current execution path and produces a section.

\subsection{Rule S4: If/Elif/Else}
\[
  \mathsf{exec}(\mathsf{if}\; c: B_1\; \mathsf{else}: B_2,
    \mathsf{env}, \mathsf{guard})
\]
\begin{enumerate}
  \item Compile the condition: $\mathsf{cond} = \mathsf{truthy}(\den{c})$.
  \item True branch: $\mathsf{guard}_T = \mathsf{guard} \land
    \mathsf{cond}$, compile $B_1$ with $\mathsf{env}_T = \mathsf{env.copy()}$.
  \item False branch: $\mathsf{guard}_F = \mathsf{guard} \land
    \lnot\mathsf{cond}$, compile $B_2$ with $\mathsf{env}_F = \mathsf{env.copy()}$.
  \item If both branches return: collect sections from both.
  \item If only one returns: collect its sections, continue with the
    other branch's modified environment.
  \item If neither returns: merge environments with
    $\mathsf{env}[x] = \mathsf{If}(\mathsf{cond}, \mathsf{env}_T[x],
    \mathsf{env}_F[x])$ for each modified variable $x$.
\end{enumerate}

\subsection{Rule S5: For Loop}
\[
  \mathsf{exec}(\mathsf{for}\; i\; \mathsf{in}\; \mathsf{range}(
    \mathsf{start}, \mathsf{stop}): B, \mathsf{env})
\]

\textbf{Case 1: Canonical fold.}  If $B$ is a single augmented
assignment $\mathsf{acc} \mathrel{\oplus}= i$ (the loop variable):
\[
  \mathsf{env}[\mathsf{acc}] \leftarrow \fold(\llbracket \oplus
    \rrbracket, \mathsf{start}, \mathsf{stop}, \mathsf{env}[\mathsf{acc}])
\]

\textbf{Case 2: General RecFunction.}  For complex loop bodies,
create a RecFunction for each modified variable:
\[
  \mathsf{lp\_acc}(i) = \begin{cases}
    \mathsf{env}[\mathsf{acc}] & \text{if } i \geq \mathsf{stop} \\
    \mathsf{body}(\mathsf{lp\_acc}(i+1), i) & \text{otherwise}
  \end{cases}
\]
Set $\mathsf{env}[\mathsf{acc}] \leftarrow \mathsf{lp\_acc}(\mathsf{start})$.

\subsection{Rule S6: While Loop}

Similar to S5 but with a guard condition instead of a range.
Create a RecFunction with a bounded counter:
\[
  \mathsf{wh\_acc}(k) = \begin{cases}
    \mathsf{env}[\mathsf{acc}] & \text{if } k > 50 \lor
      \lnot\mathsf{guard} \\
    \mathsf{body}(\mathsf{wh\_acc}(k+1)) & \text{otherwise}
  \end{cases}
\]

\subsection{Rule S7: Function Definition}
\[
  \mathsf{exec}(\mathsf{def}\; f(\ldots): B, \mathsf{env}) =
    \mathsf{env}[f \mapsto (\mathsf{\_\_func\_\_}, \mathsf{AST}(f))]
\]
Store the AST node for later inlining (D2 defunctionalization).

\subsection{Rule S8: Expression Statement (Mutation)}

For $\mathsf{obj.method}(\mathsf{args})$:
\[
  \mathsf{env}[\mathsf{obj}] \leftarrow T.\mathsf{shared\_fn}(
    \mathsf{mut\_method}, |\mathsf{args}|+1)(\mathsf{env}[\mathsf{obj}],
    \den{\mathsf{args}})
\]
This models mutation: calling a method on an object produces a new
version of that object.

\subsection{Rule S9: Try/Except}

Compile the \texttt{try} body.  For each section produced, wrap with
an uninterpreted ``handler context'' function keyed by the exception
types:
\[
  \mathsf{Section}(\mathsf{guard}, \mathsf{try\_ctx}(\mathsf{term}))
\]
where $\mathsf{try\_ctx}$ is a shared function based on the handler
signature.

\section{The Section Extractor}

The section extractor traverses a statement list and collects all
$\mathsf{Section}(\mathsf{guard}, \mathsf{term})$ pairs reachable
by any execution path:

\begin{lstlisting}[caption={Section extraction algorithm}]
def extract_sections(stmts, env, guard=True):
    sections = []
    for stmt in stmts:
        if isinstance(stmt, Return):
            sections.append(Section(guard, compile(stmt.value, env)))
            return sections  # path terminates
        elif isinstance(stmt, If):
            cond = truthy(compile(stmt.test, env))
            t_guard = And(guard, cond)
            f_guard = And(guard, Not(cond))
            # Recursively extract from both branches
            t_secs = extract_sections(stmt.body, env.copy(), t_guard)
            f_secs = extract_sections(stmt.orelse, env.copy(), f_guard)
            sections.extend(t_secs)
            sections.extend(f_secs)
            if has_return(stmt.body) and has_return(stmt.orelse):
                return sections  # both branches return
            # ... handle partial returns by continuing with modified env
        elif isinstance(stmt, For):
            exec_for(stmt, env)  # modifies env in place
        else:
            exec_one(stmt, env)  # modifies env in place
    return sections
\end{lstlisting}

The output is a list of sections forming the \emph{value presheaf}
of the program.  Each section is a local behavior under a specific
path condition (guard).


\chapter{Duck Type Inference in Detail}
\label{ch:duck-detail}

\section{The Duck Type Lattice}

Duck type inference determines which type fibers are relevant for
each parameter.  The result is a \emph{lattice element} from the
duck type lattice:

\begin{center}
\begin{tabular}{l l l}
  \textbf{Duck type} & \textbf{Allowed fibers} & \textbf{Triggering operations} \\
  \hline
  $\mathsf{int}$ & $\{\mathsf{TInt}\}$ &
    \texttt{-}, \texttt{*}, \texttt{//}, \texttt{\%}, \texttt{range(x)} \\
  $\mathsf{str}$ & $\{\mathsf{TStr}\}$ &
    \texttt{.upper()}, \texttt{.split()}, \texttt{.strip()} \\
  $\mathsf{bool}$ & $\{\mathsf{TBool}\}$ &
    (rare --- usually widened to int) \\
  $\mathsf{ref}$ & $\{\mathsf{TRef}\}$ &
    \texttt{.get()}, \texttt{.keys()}, \texttt{.items()} \\
  $\mathsf{list}$ & $\{\mathsf{TPair}, \mathsf{TRef}\}$ &
    \texttt{.append()}, \texttt{.sort()}, \texttt{.pop()} \\
  $\mathsf{collection}$ & $\{\mathsf{TPair}, \mathsf{TRef}, \mathsf{TStr}\}$ &
    \texttt{for x in p}, \texttt{len(p)}, \texttt{sorted(p)} \\
  $\mathsf{any}$ & $\{\mathsf{TInt}, \ldots, \mathsf{TRef}\}$ &
    no type-constraining operations \\
\end{tabular}
\end{center}

\section{The Inference Algorithm}

\begin{lstlisting}[caption={Duck type inference}]
def infer_duck_type(func_f, func_g, param_name):
    ops_f = extract_operations(func_f, param_name)
    ops_g = extract_operations(func_g, param_name)
    ops = ops_f | ops_g  # union of operations from both programs

    if ops & NUMERIC_ONLY_OPS:
        return 'int'    # arithmetic forces integer
    if ops & STRING_METHOD_OPS:
        return 'str'    # string methods force string
    if ops & DICT_METHOD_OPS:
        return 'ref'    # dict methods force reference
    if ops & LIST_METHOD_OPS:
        return 'list'   # list methods force list/ref
    if ops & ITERATION_OPS:
        return 'collection'  # iteration allows pair/ref/str
    return 'any'        # no constraints
\end{lstlisting}

\textbf{Key insight}: we analyze \emph{both} programs to determine
the duck type.  If program A uses arithmetic on parameter $x$ but
program B only uses comparison, the union of operations includes
arithmetic, so $x$ is constrained to $\mathsf{int}$.  This ensures
both programs are compared on the \emph{same} type fiber.

\section{Operation Extraction from AST}

The operation extractor walks the AST and collects all operations
applied to a given parameter:

\begin{lstlisting}[caption={Operation extraction}]
def extract_operations(func_node, param_name):
    ops = set()
    for node in ast.walk(func_node):
        # Binary op with param on either side
        if isinstance(node, BinOp):
            if mentions(node.left, param_name) or \
               mentions(node.right, param_name):
                ops.add(binop_name(node.op))  # 'add', 'mul', etc.

        # Method call on param
        if isinstance(node, Call) and \
           isinstance(node.func, Attribute) and \
           isinstance(node.func.value, Name) and \
           node.func.value.id == param_name:
            ops.add('method_' + node.func.attr)

        # Iteration over param
        if isinstance(node, For) and \
           isinstance(node.iter, Name) and \
           node.iter.id == param_name:
            ops.add('iter')

        # Builtin call with param as argument
        if isinstance(node, Call) and \
           isinstance(node.func, Name):
            for arg in node.args:
                if mentions(arg, param_name):
                    ops.add('call_' + node.func.id)
    return ops
\end{lstlisting}

\section{Fiber Narrowing and the Site Category}

Given duck types for all parameters, the site category is constructed:

\begin{lstlisting}[caption={Site category construction}]
def build_site_category(param_fibers):
    """param_fibers: list of lists of allowed tags per param."""
    import itertools
    combos = list(itertools.product(*param_fibers))
    # Limit combinatorial explosion
    if len(combos) > 50:
        combos = combos[:50]

    sites = {}
    for combo in combos:
        site_id = SiteId(name='fiber_' + 'x'.join(combo))
        sites[combo] = site_id

    # Overlaps: two combos that share at least one fiber component
    overlaps = []
    for i, ci in enumerate(combos):
        for j, cj in enumerate(combos):
            if i < j and any(ci[k] == cj[k] for k in range(len(ci))):
                overlaps.append((sites[ci], sites[cj]))

    return sites, overlaps
\end{lstlisting}


\chapter{The \v{C}ech Complex Construction Algorithm}
\label{ch:cech-algorithm}

This chapter gives the complete algorithm for building the \v{C}ech
complex from local equivalence judgments, computing the coboundary
operators, and determining $\Hone$.

\section{Step 1: Local Equivalence Judgments}

For each site $U_i$ (type fiber combination), run the Z3 query:

\begin{lstlisting}[caption={Local equivalence check}]
def check_local(fiber_combo, secs_f, secs_g, params, T, timeout):
    solver = Solver()
    solver.set('timeout', timeout)

    # Add semantic axioms
    T.semantic_axioms(solver)

    # Constrain params to this fiber
    for idx, fiber in enumerate(fiber_combo):
        solver.add(T.tag(params[idx]) == tag_constant(fiber))
        solver.add(T.tag(params[idx]) != T.TBot)

    # Check each section pair
    for sf in secs_f:
        for sg in secs_g:
            overlap = And(sf.guard, sg.guard)
            solver.push()
            solver.add(overlap)
            solver.add(sf.term != sg.term)
            result = solver.check()
            solver.pop()
            if result == sat:
                return INEQUIVALENT, solver.model()
            elif result == unsat:
                continue  # this pair agrees
    return EQUIVALENT, None
\end{lstlisting}

\section{Step 2: Building $C^0$}

The 0-cochain group is a vector over $\mathrm{GF}(2)$, one entry per
site:
\[
  C^0 = \mathrm{GF}(2)^{|I|}
\]
where entry $i$ is 1 if site $U_i$ is locally equivalent, 0 otherwise.

\begin{lstlisting}[caption={Building $C^0$}]
def build_c0(local_judgments):
    """Returns dict: site_id -> CochainElement."""
    c0 = {}
    for site_id, judgment in local_judgments.items():
        is_iso = (judgment.verdict == EQUIVALENT)
        c0[site_id] = CochainElement(
            sites=(site_id,),
            is_iso=is_iso,
            predicate=TruePred() if is_iso else FalsePred()
        )
    return CochainGroup(degree=0, elements=c0)
\end{lstlisting}

\section{Step 3: Computing $\delta^0 : C^0 \to C^1$}

The coboundary operator $\delta^0$ computes transition functions
on pairwise overlaps:

\begin{lstlisting}[caption={Coboundary $\delta^0$}]
def delta_0(c0, overlaps):
    """Compute C^1 from C^0 and overlap list."""
    c1 = {}
    for (site_i, site_j) in overlaps:
        phi_i = c0.at(site_i)
        phi_j = c0.at(site_j)
        if phi_i is None or phi_j is None:
            continue
        # Transition: phi_j|_{ij} o (phi_i|_{ij})^{-1}
        # Both are truth values, so transition is trivial
        # iff both are isomorphisms
        is_trivial = phi_i.is_iso and phi_j.is_iso
        c1[(site_i, site_j)] = CochainElement(
            sites=(site_i, site_j),
            is_iso=is_trivial,
            predicate=TruePred() if is_trivial else FalsePred()
        )
    return CochainGroup(degree=1, elements=c1)
\end{lstlisting}

\section{Step 4: Computing $\delta^1 : C^1 \to C^2$}

The second coboundary checks the cocycle condition on triple overlaps:

\begin{lstlisting}[caption={Coboundary $\delta^1$}]
def delta_1(c1, triple_overlaps):
    """Compute C^2 from C^1 and triple overlap list."""
    c2 = {}
    for (si, sj, sk) in triple_overlaps:
        psi_ij = c1.at(si, sj)
        psi_jk = c1.at(sj, sk)
        psi_ik = c1.at(si, sk)
        # Cocycle condition: psi_jk o psi_ij^{-1} o psi_ik
        # Over GF(2): sum mod 2
        vals = [psi_ij, psi_jk, psi_ik]
        non_trivial_count = sum(1 for v in vals
                                if v and not v.is_iso)
        is_cocycle = (non_trivial_count % 2 == 0)
        c2[(si, sj, sk)] = CochainElement(
            sites=(si, sj, sk),
            is_iso=is_cocycle
        )
    return CochainGroup(degree=2, elements=c2)
\end{lstlisting}

\section{Step 5: Computing $\Hone$}

\begin{lstlisting}[caption={$\Hone$ computation}]
def compute_h1(c0, c1, c2, overlaps, triple_overlaps):
    """Compute H^1 = ker(delta^1) / im(delta^0)."""
    # Build coboundary matrices over GF(2)
    # delta^0 matrix: rows = overlaps, cols = sites
    # delta^1 matrix: rows = triple overlaps, cols = overlaps

    n_sites = len(c0.elements)
    n_overlaps = len(overlaps)
    n_triples = len(triple_overlaps)

    # Compute rank of kernel and image
    rank_ker_d1 = n_overlaps - gf2_rank(delta1_matrix)
    rank_im_d0 = gf2_rank(delta0_matrix)

    h1_rank = rank_ker_d1 - rank_im_d0
    return h1_rank  # 0 means equivalence glues
\end{lstlisting}

\section{The GF(2) Rank Algorithm}

\begin{lstlisting}[caption={Gaussian elimination over GF(2)}]
def gf2_rank(matrix):
    """Row-reduce a binary matrix and return its rank."""
    m = [row[:] for row in matrix]  # copy
    rows, cols = len(m), len(m[0]) if m else 0
    rank = 0
    for col in range(cols):
        # Find pivot
        pivot = None
        for row in range(rank, rows):
            if m[row][col] == 1:
                pivot = row
                break
        if pivot is None:
            continue
        # Swap pivot row to position 'rank'
        m[rank], m[pivot] = m[pivot], m[rank]
        # Eliminate column in other rows
        for row in range(rows):
            if row != rank and m[row][col] == 1:
                m[row] = [(m[row][j] + m[rank][j]) % 2
                          for j in range(cols)]
        rank += 1
    return rank
\end{lstlisting}

\section{Correctness Theorem}

\begin{theorem}[Correctness of the $\Hone$ Algorithm]
The algorithm computes the rank of $\Hone(\covers, \Iso) =
\ker(\delta^1) / \mathrm{im}(\delta^0)$ over $\mathrm{GF}(2)$
correctly.  Specifically:
\begin{enumerate}
  \item $\mathrm{rank}(\Hone) = 0$ iff all 1-cocycles are
    coboundaries iff local equivalences glue.
  \item $\mathrm{rank}(\Hone) > 0$ gives the number of independent
    obstructions.
  \item The algorithm runs in $O(n^3)$ where $n$ is the number of
    overlaps.
\end{enumerate}
\end{theorem}

\begin{proof}
Standard linear algebra over a finite field.  The rank is the
dimension of the quotient space, which equals the dimension of the
kernel minus the dimension of the image.  Gaussian elimination
computes both dimensions in $O(n^3)$.
\end{proof}
\label{ch:z3-theories}

\section{The UserPropagator API}

Z3's \texttt{UserPropagator} allows implementing custom theory
solvers that participate in Z3's DPLL(T) solving loop:

\begin{lstlisting}[caption={UserPropagator skeleton}]
class MyProp(UserPropagateBase):
    def __init__(self, s=None, ctx=None):
        super().__init__(s, ctx)
        self.add_fixed(self._on_fixed)
        self.add_final(self._on_final)
        self.add_eq(self._on_eq)
        self.add_created(self._on_created)

    def push(self): ...   # save state
    def pop(self, n): ... # restore state
    def fresh(self, ctx): return MyProp(ctx=ctx)

    def _on_fixed(self, t, v): ...
    def _on_final(self): ...
    def _on_eq(self, a, b): ...
    def _on_created(self, t):
        self.add(t)  # register for tracking
\end{lstlisting}

\section{PropagateFunction}

A \texttt{PropagateFunction} is a Z3 function declaration that triggers
the propagator's \texttt{\_on\_created} callback whenever a term using
that function appears:

\begin{lstlisting}
sem = PropagateFunction('sem', IntSort(), IntSort())
# When sem(x) appears, the propagator is notified
\end{lstlisting}

\section{Planned Propagators}

\begin{enumerate}
  \item \textbf{Commutativity Propagator}: When Z3 assigns
    $\mathsf{op}(a, b) = v$, propagate $\mathsf{op}(b, a) = v$
    for commutative operations.
  \item \textbf{Fold Propagator}: Track fold terms and enforce
    permutation invariance for commutative+associative operations.
  \item \textbf{Effect Propagator}: Distinguish mutation operations
    from pure operations to prevent false equivalences.
\end{enumerate}


% ══════════════════════════════════════════════════════════
\part{Practical Engineering}
% ══════════════════════════════════════════════════════════

\chapter{End-to-End Pipeline Trace}
\label{ch:e2e}

This chapter traces a \emph{complete} equivalence check from raw
Python source to final verdict, showing every intermediate step.

\section{The Input Pair}

\begin{lstlisting}[caption={Program A: recursive sum}]
def f(n):
    if n <= 0: return 0
    return n + f(n - 1)
\end{lstlisting}

\begin{lstlisting}[caption={Program B: iterative sum}]
def g(n):
    total = 0
    for i in range(1, n + 1):
        total += i
    return total
\end{lstlisting}

\section{Step 1: Parse to AST}

Both sources are parsed by Python's \texttt{ast.parse}.  Program A
produces a \texttt{FunctionDef} with an \texttt{If} and a
\texttt{Return} containing a \texttt{BinOp(Add)} with a recursive
\texttt{Call}.  Program B produces a \texttt{FunctionDef} with an
\texttt{Assign}, a \texttt{For} over \texttt{range}, and a
\texttt{Return}.

\section{Step 2: Create Theory}

A fresh \texttt{Theory} is instantiated:
\begin{itemize}
  \item $\PyObj$ ADT created (7 constructors).
  \item $\mathsf{binop}$, $\mathsf{unop}$, $\mathsf{truthy}$,
    $\fold$ RecFunctions defined.
  \item $\mathsf{TypeTag}$ EnumSort and $\mathsf{tag}$ RecFunction
    defined.
  \item Shared function registry initialized (empty).
\end{itemize}

\section{Step 3: Compile Program A}

Parameters: $[\texttt{n}]$.  Create $p_0 = \mathsf{Const}(\texttt{p0},
\PyObj)$.  Env: $\{\texttt{n} \mapsto p_0\}$.

Detect self-recursion ($\texttt{f}$ appears in body).  Invoke
Nat-eliminator extractor:
\begin{itemize}
  \item Base guard: \texttt{n <= 0} $\Rightarrow$ threshold $= 0$.
  \item Base value: \texttt{0} $\Rightarrow$ $\mathsf{IntObj}(0)$.
  \item Step: \texttt{n + f(n-1)} $\Rightarrow$ $\op = \mathsf{ADD}$,
    param left, recursive call right.
\end{itemize}

Result: $\den{f} = \fold(\mathsf{ADD}, 0, \mathsf{ival}(p_0) + 1,
\mathsf{IntObj}(0))$.

Sections: $[(\top, \fold(\mathsf{ADD}, 0, \mathsf{ival}(p_0)+1,
\mathsf{IntObj}(0)))]$.

\section{Step 4: Compile Program B}

Parameters: $[\texttt{n}]$.  Create $p_0$ (same as A).

\texttt{total = 0}: env $\leftarrow \{\texttt{n} \mapsto p_0,
\texttt{total} \mapsto \mathsf{IntObj}(0)\}$.

\texttt{for i in range(1, n+1)}: extract range: $\mathsf{start} = 1$,
$\mathsf{stop} = \mathsf{ival}(p_0) + 1$.

Loop body: \texttt{total += i} is \texttt{AugAssign(Add)}, target
\texttt{total}, value \texttt{i} (the loop variable).  This matches
the canonical fold pattern: single accumulator, augmented by loop var.

Result: env[\texttt{total}] $\leftarrow \fold(\mathsf{ADD}, 1,
\mathsf{ival}(p_0)+1, \mathsf{IntObj}(0))$.

\texttt{return total}: section $(\top, \fold(\mathsf{ADD}, 1,
\mathsf{ival}(p_0)+1, \mathsf{IntObj}(0)))$.

\textbf{Note}: A starts fold at $0$, B starts at $1$.  But
$\fold(\mathsf{ADD}, 0, n+1, \mathsf{IntObj}(0))$ includes
$\mathsf{binop}(\mathsf{ADD}, \mathsf{IntObj}(0), \ldots) =
\ldots$ (adding 0 is identity by R3).  So both compute
$0 + 1 + 2 + \ldots + n$.

\section{Step 5: Unify Parameters}

Both have 1 parameter.  $p_0$ is shared (same Z3 Const).
No substitution needed.

\section{Step 6: Duck Type Inference}

Analyze both programs for operations on \texttt{n}:
\begin{itemize}
  \item A: comparison (\texttt{n <= 0}), subtraction (\texttt{n-1}),
    addition (\texttt{n + f(n-1)}).
  \item B: \texttt{range(1, n+1)} (addition).
\end{itemize}

Union of operations: $\{\mathsf{le}, \mathsf{sub}, \mathsf{add},
\mathsf{call\_range}\}$.  Numeric operations present $\Rightarrow$
duck type = \texttt{int}.

\section{Step 7: Build Site Category}

Single parameter, single type (int).  Site category:
\begin{itemize}
  \item One site: $U_\mathsf{int} = (\mathsf{TInt})$.
  \item No overlaps (single site).
\end{itemize}

\section{Step 8: Local Equivalence Check}

Solver created.  Semantic axioms added (none relevant --- no shared
builtins used).  Tag constraint: $\mathsf{tag}(p_0) = \mathsf{TInt}$,
$\mathsf{tag}(p_0) \neq \mathsf{TBot}$.

Section comparison: A has 1 section, B has 1 section.
Overlap: $\top \land \top = \top$ (always overlapping).

Z3 query: $\fold(\mathsf{ADD}, 0, \mathsf{ival}(p_0)+1,
\mathsf{IntObj}(0)) \neq \fold(\mathsf{ADD}, 1,
\mathsf{ival}(p_0)+1, \mathsf{IntObj}(0))$?

Z3 unfolds the fold RecFunction: the first fold includes $i=0$
(adding $\mathsf{IntObj}(0)$, which is identity by Z3's integer
arithmetic).  Both produce the same sum.

Result: \textbf{UNSAT}.  Local equivalence holds.

\section{Step 9: \v{C}ech Cohomology}

Single site, no overlaps.  $C^0 = \{1\}$ (one local isomorphism).
$C^1 = \emptyset$ (no overlaps).  $\Hone = 0$ trivially.

\section{Step 10: Verdict}

$\Hone = 0$, $\Hzero = 1$.

\textbf{EQUIVALENT}: $f \equiv g$.  H$^0$ = 1 section verified
across 1 type fiber.  The cubical path is:
$\den{f} \xrightarrow{\mathsf{D1}} \fold(\mathsf{ADD}, \ldots)
= \den{g}$.


